\documentclass[a4paper,10pt,twocolumn]{ltjarticle}
\usepackage{kaitoushu}

% グラフ用の共通スタイル定義
\pgfplotsset{
  myaxis/.style={
    axis lines=middle,
    xlabel={$x$}, ylabel={$y$},
    every axis x label/.style={at={(ticklabel* cs:1)},anchor=west},
    every axis y label/.style={at={(ticklabel* cs:1)},anchor=south},
    tick label style={font=\scriptsize},
    label style={font=\small},
    width=5.5cm, height=5cm,
    grid=none,
    samples=100,
  }
}

\begin{document}

\problemrange{3章：関数とグラフ — 練習問題1・A (p.86)}



\mondai{1}{次の関数の（\quad）内の定義域に対する値域を求めよ．}
\kotae{
$y=2x^{2}$ は下に凸の放物線．頂点は $(0,0)$．

(1) $y=2x^{2}\;(1\leqq x<2)$

区間 $[1,2)$ では単調増加．
$x=1$：$y=2$，$x=2$：$y=8$（含まない）

$\therefore 2\leqq y<8$

(2) $y=2x^{2}\;(-1\leqq x<2)$

頂点 $x=0$ は定義域内．
$x=0$：$y=0$（最小），$x=-1$：$y=2$，$x=2$：$y=8$（含まない）

$\therefore 0\leqq y<8$
}

\mondai{2}{次の条件を満たし，$y$軸に平行な軸をもつ放物線の方程式を求めよ．}
\kotae{
(1) 頂点 $(2,4)$，点 $(0,1)$ を通る

$y=a(x-2)^{2}+4$．$(0,1)$ 代入：$4a+4=1$，$a=-\dfrac{3}{4}$

$\therefore y=-\dfrac{3}{4}(x-2)^{2}+4$

(2) 軸 $x=1$，2点 $(0,-1)$，$(3,-10)$ を通る

$y=a(x-1)^{2}+q$．
$(0,-1)$：$a+q=-1$\quad $(3,-10)$：$4a+q=-10$

辺々引いて $3a=-9$，$a=-3$，$q=2$

$\therefore y=-3(x-1)^{2}+2$

(3) 3点 $(0,4)$，$(3,1)$，$(4,-4)$ を通る

$y=ax^{2}+bx+c$．$c=4$．
$(3,1)$：$9a+3b=-3$，$3a+b=-1$
$(4,-4)$：$16a+4b=-8$，$4a+b=-2$

差より $a=-1$，$b=2$

$\therefore y=-x^{2}+2x+4$
}

\mondai{3}{次の2次関数の最大値または最小値を求めよ．}
\kotae{
平方完成して頂点を読む．$a>0$なら頂点で最小，$a<0$なら頂点で最大．

(1) $y=x(x-2)=x^{2}-2x=(x-1)^{2}-1$

$x=1$ で最小値 $-1$（最大値なし）

(2) $y=2x^{2}+3x+2=2\!\left(x+\dfrac{3}{4}\right)^{2}+\dfrac{7}{8}$

$x=-\dfrac{3}{4}$ で最小値 $\dfrac{7}{8}$（最大値なし）

(3) $y=-\dfrac{1}{2}x^{2}-5x+1=-\dfrac{1}{2}(x+5)^{2}+\dfrac{27}{2}$

$x=-5$ で最大値 $\dfrac{27}{2}$（最小値なし）
}

\mondai{4}{次の関数の最大値および最小値を求めよ．ただし，（\quad）内は定義域である．}
\kotae{
手順：平方完成 → 頂点が定義域内か判定 → 端点と頂点の値を比較．

(1) $y=x^{2}-4x+3=(x-2)^{2}-1\;(-1\leqq x\leqq 5)$

頂点 $x=2$ は定義域内．
$x=-1$：$8$，$x=2$：$-1$，$x=5$：$8$

最大値 $8$（$x=-1,5$），最小値 $-1$（$x=2$）

(2) $y=-3x^{2}-6x+5=-3(x+1)^{2}+8\;(-4\leqq x\leqq -1)$

頂点 $x=-1$ は定義域の端．
$x=-4$：$-3\cdot9+8=-19$，$x=-1$：$8$

最大値 $8$（$x=-1$），最小値 $-19$（$x=-4$）

(3) $y=-x^{2}+3x-\dfrac{1}{4}=-\!\left(x-\dfrac{3}{2}\right)^{2}+2\;(1\leqq x\leqq 3)$

頂点 $x=\dfrac{3}{2}$ は定義域内．
$x=1$：$\dfrac{7}{4}$，$x=\dfrac{3}{2}$：$2$，$x=3$：$-\dfrac{1}{4}$

最大値 $2$（$x=\dfrac{3}{2}$），最小値 $-\dfrac{1}{4}$（$x=3$）
}

\mondai{5}{次の2次不等式を解け．}
\kotae{
(1) $6x^{2}-13x+6\leqq0$

因数分解：$(2x-3)(3x-2)\leqq0$

根：$x=\dfrac{3}{2},\dfrac{2}{3}$．$a>0$なので根の間：

$\therefore \dfrac{2}{3}\leqq x\leqq\dfrac{3}{2}$

(2) $2x^{2}-x+3<0$

$D=1-24=-23<0$．$a>0$より常に正．

$\therefore$ 解なし
}

\mondai{6}{2次関数 $y=x^{2}-3x-4$ のグラフを $x$ 軸方向にどれだけ平行移動すれば，原点を通るようにできるか．}
\kotae{
$x$軸方向に $p$ 平行移動すると $y=(x-p)^{2}-3(x-p)-4$．

原点を通る条件：$x=0,\;y=0$ を代入
\[
p^{2}+3p-4=0 \;\Longrightarrow\; (p+4)(p-1)=0
\]

$\therefore p=-4$ または $p=1$

（$x$ 軸方向に $-4$ または $1$ 平行移動すればよい）
}

\mondai{7}{放物線 $y=x^{2}-2x+2$ は，放物線 $y=x^{2}+4x+3$ をどのように平行移動したものか．}
\kotae{
両方を標準形に直して頂点を比較．

$y=x^{2}-2x+2=(x-1)^{2}+1$\quad 頂点 $(1,1)$

$y=x^{2}+4x+3=(x+2)^{2}-1$\quad 頂点 $(-2,-1)$

頂点の移動：$(-2,-1)\to(1,1)$

$\therefore x$軸方向に $3$，$y$軸方向に $2$ 平行移動したもの
}

\mondai{8}{放物線 $y=x^{2}+bx+c$ を $x$ 軸方向に $-2$，$y$ 軸方向に $5$ 平行移動した放物線の頂点が $(0,2)$ であるとき，定数 $b$，$c$ の値を定めよ．}
\kotae{
元の放物線：$y=x^{2}+bx+c=\!\left(x+\dfrac{b}{2}\right)^{2}-\dfrac{b^{2}}{4}+c$

頂点は $\left(-\dfrac{b}{2},\;c-\dfrac{b^{2}}{4}\right)$．

平行移動後の頂点：$\left(-\dfrac{b}{2}-2,\;c-\dfrac{b^{2}}{4}+5\right)=(0,2)$

$x$座標より：$-\dfrac{b}{2}-2=0 \Rightarrow b=-4$

$y$座標より：$c-\dfrac{16}{4}+5=2 \Rightarrow c=1$

$\therefore b=-4,\;c=1$
}

\mondai{9}{$k$ が定数のとき，2次方程式 $x^{2}-4x+k+5=0$ の実数解の個数を調べよ．}
\kotae{
判別式
\[
\dfrac{D}{4}=4-(k+5)=-k-1
\]

\begin{itemize}
\item $\dfrac{D}{4}>0 \Leftrightarrow k<-1$：異なる2つの実数解（2個）
\item $\dfrac{D}{4}=0 \Leftrightarrow k=-1$：重解（1個）
\item $\dfrac{D}{4}<0 \Leftrightarrow k>-1$：実数解なし（0個）
\end{itemize}
}

\end{document}
